\documentclass[11pt]{article}
\usepackage[a4paper,margin=1in]{geometry}
\usepackage{fontspec}
\usepackage[boldfont=false]{xeCJK}
\setCJKmainfont{FandolSong-Regular.otf}
\usepackage{amsmath}
\usepackage{amssymb}
\usepackage{array}
\usepackage{booktabs}
\usepackage{graphicx}
\usepackage{adjustbox}
\usepackage{enumitem}
\setlist[itemize]{topsep=2pt,partopsep=0pt,parsep=0pt,itemsep=2pt,leftmargin=1.4em}
\setlist[enumerate]{topsep=2pt,partopsep=0pt,parsep=0pt,itemsep=2pt,leftmargin=1.6em}
\usepackage{parskip}
\usepackage{fvextra}
\fvset{breaklines=true,breakanywhere=true,fontsize=\small}
\setlength{\parindent}{0pt}

\begin{document}

\begin{center}
  {\LARGE\bfseries Nitrogen compounds}\\[0.35em]
  {\large A Level Chemistry}
\end{center}
\vspace{0.6em}

\section*{Primary amines}

An \textbf{amine} 胺 has an $\text{–NH}_2$ group (the nitrogen replaces a hydrogen of ammonia).

\subsection*{Making a primary amine}

Heat a \textbf{halogenoalkane} 卤代烷 with ammonia dissolved in ethanol, under pressure:

\[
\text{C}_2\text{H}_5\text{Br} + \text{NH}_3 \rightarrow \text{C}_2\text{H}_5\text{NH}_2 + \text{HBr}
\]

This is a \textbf{nucleophilic substitution} 亲核取代: the lone pair on the nitrogen of ammonia attacks the slightly positive carbon and pushes out the halogen. You use an excess of ammonia, or the amine made can react again.

\section*{Nitriles and hydroxynitriles}

\subsection*{Making a nitrile}

Heat a halogenoalkane with potassium cyanide ($\text{KCN}$) in ethanol:

\[
\text{C}_2\text{H}_5\text{Br} + \text{KCN} \rightarrow \text{C}_2\text{H}_5\text{CN} + \text{KBr}
\]

This is also a nucleophilic substitution, with the $\text{CN}^-$ ion as the nucleophile. It is useful because it \textbf{adds one carbon} to the chain. The product is a \textbf{nitrile} 腈.

\subsection*{Making a hydroxynitrile}

Add $\text{HCN}$ (with $\text{KCN}$ as catalyst, and heat) to an \textbf{aldehyde} 醛 or \textbf{ketone} 酮. The $\text{H}$ and $\text{CN}$ add across the C=O bond to give a \textbf{hydroxynitrile} 羟基腈. The reagent is \textbf{hydrogen cyanide} 氰化氢, and the mechanism is \textbf{nucleophilic addition} 亲核加成.

\subsection*{Hydrolysis of nitriles}

Warm a nitrile with dilute acid (or dilute alkali, then acidify). This \textbf{hydrolysis} 水解 turns the $\text{–CN}$ group into a $\text{–COOH}$ group, giving a \textbf{carboxylic acid} 羧酸:

\[
\text{CH}_3\text{CN} + 2\text{H}_2\text{O} + \text{HCl} \rightarrow \text{CH}_3\text{COOH} + \text{NH}_4\text{Cl}
\]

A nitrile can also be \textbf{reduced} by hydrogen and a catalyst to form an amine, which is the \textbf{reduction} 还原 route to a longer-chain amine.


\end{document}
